\chapter*{Preface}
\addcontentsline{toc}{chapter}{Preface}
\label{ch:preface}

\section*{About This Manual}

This manual documents SIRCUS, the SIRC User System: the console built from the SIRC-1 CPU, the SIRCADA display adapter
and the SIRCAS sound processor. The three chips each have a manual of their own, which describes what the chip does and
how it behaves at its pins. This manual describes what the chips do not describe for themselves: how they are
connected, which clock each one runs from, how they share one bus and a small set of interrupt lines, how they are
reset, and how the timing of each one relates to the others.

It is a companion to the \archmanual{}, which defines the CPU's instructions and exceptions, to the \hwmanual{}, which
defines the CPU package and bus, to the \asmmanual{}, which defines the language used to program the system, to the
\videomanual{} and to the \audiomanual{}.

It is written for:

\begin{itemize}
	\item Engineers building the console, who need to know what every pin connects to and how fast
	\item Programmers who need the timing relationships between the picture, the sound and the CPU
	\item Authors of tools and emulators, who must reproduce those relationships
\end{itemize}

\section*{Assumed Knowledge}

This manual assumes the reader has the manual of each chip to hand. It repeats a fact from them only where the fact is
needed to understand the connection, and it names the chapter that holds the full account instead of reproducing it.

\section*{Revision Status}

This is a draft of the \sysmanual{}, not for distribution. Anything not yet decided in the manual is marked ``not yet
specified''; the manual does not invent values for it.
